\chapter{Experiment Tracking with MLflow and DagsHub}\label{ch:tracking}

\section{Tool lab: MLflow locally}
\file{labs/mlflow_toy.py} trains a network with one hidden layer on 10\,000 Fashion-MNIST images.
It logs the parameters (\texttt{lr}, \texttt{epochs}, \texttt{units}, \texttt{optimizer}), the
per-epoch metrics through a small Keras callback (\texttt{mlflow.log\_metrics(..., step=epoch)}),
a loss plot (\texttt{mlflow.log\_figure}) and the model (\texttt{mlflow.tensorflow.log\_model}). I
ran it twice, with \texttt{lr=0.001} and \texttt{lr=0.01}, against a local store (\cmd{--local}
writes to \file{./mlruns_toy}), and then opened \cmd{mlflow ui --backend-store-uri ./mlruns_toy}.
The comparison view overlays the two \texttt{val\_accuracy} curves and highlights the parameter
that differs. These two runs ran on CPU in a separate Linux workspace (TensorFlow 2.20, MLflow
2.22.5) with the UI on \texttt{127.0.0.1}.

After 5 epochs the higher learning rate fits the training data faster (training accuracy 0.862
against 0.849, loss 0.373 against 0.437) but does no better on validation (0.842 against 0.846),
so I would keep \texttt{lr=0.001}.

\shot{s3_1_mlflow_runs}{Local MLflow UI: the toy experiment with both runs.}
\shot{s3_1_mlflow_compare}{Comparison of the two learning rates.}

In my own words, the MLflow terms are:
\begin{description}[leftmargin=2.2cm,style=nextline]
\item[Experiment] a named group of runs working on the same problem, so they can be compared side
by side. Here I have \texttt{VisionGuard} and \texttt{toy-fashion-mnist}.
\item[Run] one execution of the training code. It has a unique id, a start and end time and a
status, and everything logged while it runs belongs to it.
\item[Parameter] an input setting that is fixed for the run and written once, such as the learning
rate. Parameters cannot be changed: logging a different value for the same key raises an error.
\item[Metric] a numeric result that can be logged many times with a step, for example
\texttt{val\_loss} after every epoch. This is how MLflow draws the curves.
\item[Artifact] any output file stored with the run, such as plots, CSVs, the requirements list or
the model itself.
\item[Tag] free-form key/value information that describes the run without configuring it, such as
the git commit, the data version or who ran it. Tags can be edited later and are useful for
searching.
\end{description}

\section{DagsHub setup}
I created the DagsHub project \texttt{23L-0764\_A2\_Ops} with ``Connect a repository'', so it
mirrors my GitHub repository and adds DVC storage and a hosted MLflow server. The tracking address
is \url{\MLflowURL}, but it does not appear anywhere in the code. \file{tracking.py} calls
\texttt{load\_dotenv()} and reads \texttt{MLFLOW\_TRACKING\_URI}, and MLflow reads
\texttt{MLFLOW\_TRACKING\_USERNAME} and \texttt{MLFLOW\_TRACKING\_PASSWORD} (the DagsHub token) from
the environment. These values are in \file{.env}, which is listed in \file{.gitignore} and
\file{.dockerignore}, and \file{.env.example} shows the variable names with dummy values. If
\texttt{MLFLOW\_TRACKING\_URI} is not set, the code falls back to a local \file{mlruns/} folder, so
it still works offline.

\shot{s3_2_dagshub_project}{The DagsHub project connected to GitHub.}
\shot{s3_2_dagshub_toy_run}{The toy run logged to DagsHub's MLflow server.}
\shot{s3_2_gitignore}{\file{.env} is git-ignored: \cmd{git check-ignore -v .env} names the rule and \cmd{git status} does not list it.}

\section{Instrumenting VisionGuard training}
Each run of \file{scripts/train.py} logs:
\begin{itemize}
\item every value from \file{params.yaml} that was actually used, flattened into names like
\texttt{train.lr\_head}, plus the split sizes and class weights, as parameters;
\item loss, accuracy, AUC, precision and recall for training and validation after every epoch, and
the learning rate, so ReduceLROnPlateau shows up in the curves;
\item final metrics: \texttt{best\_epoch}, the chosen \texttt{val\_threshold}, all validation and
test metrics at that threshold (\texttt{test\_accuracy}, \texttt{test\_f1},
\texttt{test\_roc\_auc}, \texttt{test\_false\_positive\_rate}, the confusion-matrix cells), and
test accuracy and F1 at the plain 0.5 threshold for comparison;
\item tags: \texttt{git\_commit}, \texttt{git\_dirty}, \texttt{dvc\_data\_version} (the md5 from
\file{data/raw.dvc}), \texttt{seed}, the Python, TensorFlow and Keras versions, the platform, model
type, backbone and tracking mode.
\end{itemize}

With manual tracking (\cmd{--tracking manual}, for example the run \texttt{mnv2\_default}) my
callback \texttt{MLflowEpochLogger} passes Keras' \texttt{logs} to \texttt{mlflow.log\_metrics}
at the end of every epoch with a global step, so the two training stages form one continuous curve.
With autologging (\cmd{--tracking autolog}, run \texttt{mnv2\_autolog}),
\texttt{mlflow.tensorflow.autolog()} patches \texttt{model.fit()} and records things by itself.

\begin{table}[H]\centering\small
\begin{tabularx}{\linewidth}{l X X}
\toprule
 & Manual callback & Autologging \\
\midrule
Captures & exactly what I choose: per-epoch metrics and LR, my parameters and tags & every \texttt{fit()} argument (epochs, batch size, class weight\ldots), all optimizer settings (\texttt{opt\_*}), per-epoch metrics, model summary and EarlyStopping info, with no extra code \\
Misses & anything I forget to log & test metrics, the chosen threshold, git/DVC tags, the confusion matrix and the splits, because it only sees \texttt{fit()} \\
Two-stage training & one continuous run & the second \texttt{fit()} tries to log \texttt{epochs} again with a new value, which is not allowed, so I put stage 2 in a nested child run \\
Clutter & short & many parameters that have nothing to do with the experiment \\
\bottomrule
\end{tabularx}
\caption{Manual logging versus autologging.}\end{table}

In practice I would use both together: autolog for what happens inside training, and explicit
logging for provenance and evaluation. Both kinds of run here came from the same configuration.

\shot{s3_3_manual_run}{Manual run: parameters, metric curves and provenance tags.}
\shot{s3_3_autolog_run}{Autologged run: the automatically captured \texttt{fit()} and optimizer parameters.}

I hit two errors while setting this up. MLflow rejected my first metric name,
\texttt{test@0.5\_accuracy}, because metric names may only contain letters, digits, \texttt{\_ - .
/ :} and spaces, so I renamed it to \texttt{test\_t0.5\_accuracy}. Then the first autolog run had
no epoch metrics at all. Autolog tried to log \texttt{class\_weight}, which I had already logged
under the same key in a different string format, and because parameters cannot change, the autolog
hook gave up. I renamed my parameter to \texttt{class\_weights} and moved the second \texttt{fit()}
into a nested run.

\section{Artifacts}
Each run stores \file{plots/} (training curves, confusion matrix, ROC curve), \file{splits/} (the
exact train, validation and test CSVs), \file{environment/} (\file{requirements-freeze.txt} from
\cmd{pip freeze}, and \file{params.yaml}), \file{evaluation/test_predictions.csv},
\file{model_info/} (model summary and the trainable layers per stage), \file{keras_model/} (the
\file{model.keras} file and \file{model_meta.json} with image size, resize mode and threshold) and
\file{model/}, the MLflow TensorFlow-flavour model with an input/output signature that the registry
uses.

\shot{s3_4_artifacts}{Artifacts tab of the champion run.}

\section{Experiment comparison}\label{sec:3-5}
\file{scripts/run_experiments.py} launches the reference run and a set of variants, each changing
one thing compared with the reference, so any difference can be put down to that change.
\file{scripts/compare_runs.py} pulls all finished runs from MLflow and builds the table and the
chart.

\gentable{runs}{All VisionGuard runs (sorted by validation ROC-AUC).}
\autofig{runs_comparison}{Validation ROC-AUC and test F1 per run.}
\shot{s3_5_compare}{MLflow/DagsHub run comparison view.}

I pick the model by validation ROC-AUC, which does not depend on a threshold and measures how well
the model ranks AI images above real ones. Ties are broken by the validation false-positive rate,
since false accusations are the expensive error. The test metrics are in the table but play no part
in the choice; otherwise the test set would turn into a second validation set. The chosen run is
\texttt{\vgBestRun}, with validation AUC \vgBestValAuc, test AUC \vgBestTestAuc, test accuracy
\vgBestTestAcc{} and test F1 \vgBestTestFone.

Starting the transfer stage from the tuner's settings helped the most (validation AUC from 0.830 to
0.881, test F1 from 0.676 to 0.731). Turning augmentation off also improved the validation AUC, to
0.859. On a dataset this small the random flips, rotations and zoom mostly slowed learning down,
and zoom and rotation also change the fine generator textures the model depends on. The 160\,px
input did worst on the selection metric (validation AUC 0.764, validation FPR 0.29), since
shrinking the image wipes out fine texture. It had a higher test accuracy, but on 146 images that
is noise, and the selection rule ignores it on purpose. EfficientNetB0 had the best test AUC
(0.852) but a lower validation AUC, so I did not choose it. The autolog run matches the manual
reference (validation AUC 0.830 for both; test AUC 0.764 against 0.754, because GPU training is
not bit-exact).

My first version of the comparison ranked runs by \texttt{val\_auc}. That is the Keras AUC from the
last epoch the callback logged, which is not the best checkpoint that actually gets evaluated, and
autologged runs did not have it with the same meaning. I switched to \texttt{val\_roc\_auc}, which
my evaluation code computes on the reloaded best model for every run.

\section{Model registry}
\file{scripts/register_model.py} makes sure the registered model \texttt{VisionGuard-23L-0764} has
at least two versions. The reference transfer run and the baseline CNN are registered during
training with \cmd{--register}, and the selected run is registered if it is not there yet. The
script gives the best version the alias \texttt{champion} and the runner-up the alias
\texttt{challenger}, and adds a description and a \texttt{val\_auc} tag. An alias works like a
pointer that can be moved. To promote a new model I move the alias, and anything that loads
\texttt{@champion} picks up the new version with no code change.

\gentable{registry}{Registered versions and aliases.}
\shot{s3_6_registry}{Model registry on DagsHub with the \texttt{champion} and \texttt{challenger} aliases.}

\file{scripts/load_champion.py} is a separate script. It resolves the alias with
\texttt{MlflowClient().get\_model\_version\_by\_alias()}, loads that version with
\texttt{mlflow.tensorflow.load\_model("models:/VisionGuard-23L-0764/<v>")}, reads the threshold and
image size from the version's \file{model_meta.json} and classifies one image.

\shot{s3_6_load_alias}{Loading the model by alias and classifying one image.}

The serving processes should not contact DagsHub on every request, so
\file{scripts/export_champion.py} resolves \texttt{@champion} once, downloads \file{model.keras} and
writes \file{models/champion.keras} together with \file{champion.json} (version, threshold, image
size). It saves the model again without the optimizer state, since serving never trains, which
brings it down to 9.6\,MB. I committed this exported copy to Git. It is small, only used for
inference and can be regenerated from the registry, which remains the source of truth, and having
it in the repository lets CI and the cloud VM build the API image without registry credentials.
Locally the \file{models/} folder is mounted read-only into the web container; for the cloud API
the file is copied into the image at build time (\cref{ch:docker,ch:deploy}).
